\section{Time Management and Constraints}
\label{sec:time_management}

In practical chess engines, search is not only limited by depth but often by a
set of external constraints such as time, node budgets, or explicit stop
signals. In this engine, all stopping conditions for a single search invocation
are bundled in an immutable constraint object, which is initialized once before
the search starts and treated as read-only during the search.

A key design goal is that the search code itself does not depend on UCI parsing
details. Instead, the UCI layer produces a fully specified \texttt{SearchConstraints}
instance, and the search only consults this object (and a stop flag) to decide
when to terminate.

\subsection{Unified Constraint Model}

The engine uses a dedicated \texttt{SearchConstraints} structure to describe how
a search is controlled and when it should terminate. This structure unifies all
search modes under a single interface and avoids scattered termination logic
across the search code.

A constraint instance contains:
\begin{itemize}
    \item a \textbf{search mode} determining which criteria are active,
    \item an optional \textbf{fixed move time} (\texttt{movetime\_ms\_}),
    \item an optional \textbf{fixed depth} (\texttt{depth\_}),
    \item an optional \textbf{node limit} (\texttt{nodes\_}),
    \item and a \textbf{time budget} object (\texttt{budget\_}) used for time-managed searches.
\end{itemize}

To simplify checks in the search, convenience predicates such as
\texttt{has\_fixed\_depth()}, \texttt{has\_node\_limit()}, and \texttt{has\_movetime()}
are provided.

\subsection{Search Modes}

The active search behavior is selected via \texttt{SearchType}. The engine
supports the following modes:
\begin{itemize}
    \item \textbf{Normal:} Tournament-style time management with iterative
    deepening (default UCI behavior). Time allocation is derived from the clock
    settings and stored in \texttt{budget\_}.
    \item \textbf{FixedTime:} Search for a fixed amount of time specified by
    \texttt{movetime\_ms\_} (UCI \texttt{go movetime}).
    \item \textbf{FixedDepth:} Search to a fixed depth only, determined by
    \texttt{depth\_}, without relying on time control.
    \item \textbf{NodeLimit:} Search until a fixed number of nodes specified by
    \texttt{nodes\_} has been explored.
    \item \textbf{Infinite:} Search indefinitely until an external stop command
    is received (UCI \texttt{go infinite}).
    \item \textbf{Ponder:} Search during the opponent’s time. The engine
    continues searching until it is either stopped or the predicted opponent
    move is played, at which point the search can be continued or restarted.
\end{itemize}

This mode-based design enables a clean separation between the UCI interface
(which interprets the \texttt{go} command) and the search itself (which only
consults the active constraints).

\subsection{Time Budgets: Structure and Arming}

Time-based searches (Normal / FixedTime) rely on \texttt{TimeControl} stored in
\texttt{SearchConstraints::budget\_}. The budget distinguishes a \emph{soft} and a
\emph{hard} limit:
\begin{itemize}
    \item \textbf{Soft deadline:} preferred stopping point, checked only at safe
    boundaries (between completed iterations of iterative deepening).
    \item \textbf{Hard deadline:} strict upper bound that must not be exceeded,
    checked frequently throughout the search (including quiescence).
\end{itemize}

\paragraph{Unarmed Budgets for Non-time Modes}
A crucial implementation detail is that \texttt{budget\_} is explicitly
\emph{disarmed} for non-time-based modes. Concretely, deadlines are set to zero,
so \texttt{soft\_time\_up(...)} and \texttt{hard\_time\_up(...)} always return
false. This ensures that depth-limited and node-limited search do not depend on
any timing state.

\paragraph{Arming the Budget Before Search}
Immediately before a time-managed search begins, \texttt{TimeManager::init\_search}
is called. This function captures the current monotonic time
(\texttt{now\_ms()}) as the start timestamp and derives absolute deadlines:
\[
\texttt{soft\_deadline\_ms\_} = \texttt{start\_ms\_} + \texttt{soft\_ms\_}, \qquad
\texttt{hard\_deadline\_ms\_} = \texttt{start\_ms\_} + \texttt{hard\_ms\_}.
\]
From this point onward, the budget object is treated as read-only.

\subsection{Termination Conditions in the Search}

The engine enforces termination using a combination of (i) an external stop flag,
(ii) the hard deadline, and (iii) an optional node limit. These checks are
centralized in \texttt{ChessBot::hard\_stop()}, which is called frequently inside
the main search and quiescence.

\paragraph{Hard Stop Checks}
The stop check follows the order shown in \cref{algo:hardstop}: first an external
flag, then the hard time deadline, then the node limit. If any condition is met,
the search aborts immediately and returns a special ``aborted'' result to the
caller.

\begin{algorithm}[ht]
\DontPrintSemicolon
\KwIn{Current time $now$}
\KwOut{\texttt{true} iff the search must terminate immediately}

\If{\texttt{stop\_requested} is set}{
    \Return \texttt{true}\;
}
\If{\texttt{budget\_.hard\_time\_up(now)}}{
    \Return \texttt{true}\;
}
\If{\texttt{has\_node\_limit()} \textbf{and} \texttt{nodes >= nodes\_}}{
    \Return \texttt{true}\;
}
\Return \texttt{false}\;
\caption{Immediate termination check (\texttt{hard\_stop()}).}
\label{algo:hardstop}
\end{algorithm}

\paragraph{Soft Stop in Iterative Deepening}
In contrast to the hard deadline, the \emph{soft} deadline is checked only after
a complete iteration has finished. This ensures that the engine does not stop
mid-iteration, which would otherwise risk returning an unstable or partially
explored root result. The structure is summarized in \cref{algo:softstop}.

\begin{algorithm}[ht]
\DontPrintSemicolon
\KwIn{Board position $P$}
\KwOut{Best move found within the budget}

$bestMove \gets$ legal fallback move\;
\For{$d \gets 1,2,3,\dots$}{
    $(score, aborted) \gets$ depth-limited root search at depth $d$\;
    \If{$aborted$}{
        \textbf{break}\tcp*{do not trust partial iteration}
    }
    $bestMove \gets$ move from completed iteration\;
    \If{\texttt{budget\_.soft\_time\_up(now)}}{
        \textbf{break}\tcp*{stop cleanly after a full iteration}
    }
}
\Return $bestMove$\;
\caption{Soft stopping policy at iteration boundaries (iterative deepening).}
\label{algo:softstop}
\end{algorithm}

\subsection{Summary}

The engine models time control and other stopping criteria using a single,
immutable \texttt{SearchConstraints} object. By activating different termination
rules depending on \texttt{SearchType} and by explicitly arming or disarming time
budgets, the implementation remains robust and easy to analyze.

Most importantly, the implementation distinguishes between \emph{hard} stopping
conditions, which are enforced immediately inside the recursive search, and a
\emph{soft} stopping condition, which terminates the search cleanly between
iterations of iterative deepening. This separation is essential for making
time-constrained iterative deepening comparable to fixed-depth search in later
experiments, since it ensures that the engine always returns the best move from
the last fully completed depth.